% =============================================================================
% CHAPTER 23: MULTI-LEVEL IMPLEMENTATION
% Scaling the Framework from Individual to Global
% =============================================================================
%
% Version: 1.0 (January 2026) - New multi-level chapter
%
% Purpose: Demonstrate how Chapters 16-22 framework scales across organizational
% and political levels. Shows that 9C + Portfolio logic applies at every level
% (Individual, Household, Organization, National, Global), but parameters,
% interactions (γ), and spillovers differ systematically by level.
%
% This chapter bridges from policy theory to implementation reality by showing
% how decisions cascade through levels and how constraints differ at each level.
%
% For readers: Government implementation teams, organizational leaders,
% international coordinators who need to understand multi-level coordination.
%
% Primary Appendix: NNN (METHOD-MULTI-LEVEL)
% Secondary Appendix: HHH (toolkit), KKK (dynamics), LLL (templates), MMM (coherence)
%
% Prerequisites: Chapters 16-22 (complete portfolio framework)
% Continues from: Chapter 22 (Policy Applications)
% Leads to: Future specialized implementations
%
% Chapter Type: C (Application, Multi-Level)
%
% =============================================================================

\chapter{Multi-Level Implementation: Scaling Across Individuals, Organizations, and Nations}
\label{ch:multi-level-implementation}

%%%%%%%%%%%%%%%%%%%%%%%%%%%%%%%%%%%%%%%%%%%%%%%%%%%%%%%%%%%%%%%%%%%%%%%%%%%%%%%
% CHAPTER SCOPE BOX
%%%%%%%%%%%%%%%%%%%%%%%%%%%%%%%%%%%%%%%%%%%%%%%%%%%%%%%%%%%%%%%%%%%%%%%%%%%%%%%
\begin{tcolorbox}[
    colback=purple!5!white,
    colframe=purple!75!black,
    title={\textbf{Chapter Scope: Ziel / In-Scope / Out-of-Scope}},
    fonttitle=\bfseries
]

\textbf{Ziel (Objective):}
Apply the EBF framework to organizational scaling with practical examples and policy implications.

\vspace{0.3cm}
\textbf{In-Scope:}
\begin{itemize}[nosep]
    \item Practical applications of organizational scaling
    \item Multiple worked examples (≥2)
    \item Policy implications and recommendations
    \item Implementation guidance
\end{itemize}

\vspace{0.3cm}
\textbf{Out-of-Scope (delegiert an Appendices):}
\begin{itemize}[nosep]
    \item Theoretical foundations $\rightarrow$ Foundation chapters
    \item Formal proofs $\rightarrow$ FORMAL appendices
    \item Parameter estimation methods $\rightarrow$ METHOD appendices
\end{itemize}

\end{tcolorbox}



%%%%%%%%%%%%%%%%%%%%%%%%%%%%%%%%%%%%%%%%%%%%%%%%%%%%%%%%%%%%%%%%%%%%%%%%%%%%%%%
% CHAPTER OVERVIEW
%%%%%%%%%%%%%%%%%%%%%%%%%%%%%%%%%%%%%%%%%%%%%%%%%%%%%%%%%%%%%%%%%%%%%%%%%%%%%%%
\subsection*{Chapter Overview}

This chapter is organized as follows:

\begin{itemize}
\item \textbf{Section \ref{sec:ch23-intro}}: Introduction and motivation
\item \textbf{Section \ref{sec:ch23-theory}}: Theoretical foundations
\item \textbf{Section \ref{sec:ch23-applications}}: Applications and examples
\item \textbf{Section \ref{sec:ch23-summary}}: Summary and conclusions
\end{itemize}



% =============================================================================
% QUICK REFERENCE BOX
% =============================================================================

\begin{tcolorbox}[colback=blue!5!white, colframe=blue!75!black, title=Key Concepts in This Chapter]
\small
\begin{tabular}{@{}ll@{}}
\textbf{Organizational Level} & Single person making decisions \\
\textbf{Individual (L0)} & Single person making decisions \\
\textbf{Household (L1)} & Couple/family with shared context \\
\textbf{Organization (L2)} & Company/department with hierarchy \\
\textbf{National (L3)} & Regional/national government \\
\textbf{Global (L4)} & International coordination \\
\textbf{9C at Every Level} & Framework structure unchanged \\
\textbf{γ by Level} & Complementarities differ by level \\
\textbf{Spillovers} & Cross-level effects \\
\textbf{Aggregation} & How micro behavior becomes macro outcomes \\
\end{tabular}
\end{tcolorbox}

% =============================================================================
% APPENDIX REFERENCES
% =============================================================================

\begin{tcolorbox}[colback=yellow!10!white, colframe=orange!75!black, title=Appendix References for This Chapter]

\textbf{Essential:}
\begin{itemize}[nosep]
\item Appendix UNMAPPED_NNN (METHOD-MULTI-LEVEL) --- Templates and parameter tables for all 5 levels
\item Appendix UNMAPPED_MMM (METHOD-POLICY-COHERENCE) --- Cross-level spillover analysis
\end{itemize}

\textbf{Primary Appendix:} Appendix UNMAPPED_NNN

\textbf{Supporting:}
\begin{itemize}[nosep]
\item Appendix TKT (METHOD-TOOLKIT) --- Coherence criteria at all levels
\item Appendix UNMAPPED_KKK (METHOD-DYNAMIC) --- Formal dynamics (apply to each level)
\item Appendix UNMAPPED_LLL (METHOD-POLICY-DOMAINS) --- Domain-specific templates
\item \textbf{Appendix UNMAPPED_V (CORE-WHEN)} --- Network structure parameter $\rho$ for narrative transmission
\item \textbf{Appendix UNMAPPED_RB (LIT-BENABOU)} --- Moral narratives, network segregation effects
\end{itemize}

\textit{Note: Network structure ($\rho$) from BCM2 factor CH-SOC-13 moderates how narratives spread across organizational levels. Homophilic networks ($\rho \to 1$) amplify polarization; mixing ($\rho \to 0$) enables convergence.}

\end{tcolorbox}

% =============================================================================
% CROSS-REFERENCE: EMERGE ALGORITHM (Chapter 11, EIT-EMP)
% =============================================================================
\begin{tcolorbox}[colback=blue!5!white, colframe=blue!75!black,
    title=\textbf{Cross-Reference: EMERGE Algorithm \& Multi-Level K* (Chapter 11, 20)}]

\textbf{Zwei Level-Konzepte (nicht verwechseln!):}
\begin{center}
\small
\begin{tabular}{@{}lll@{}}
\textbf{Konzept} & \textbf{Chapter 23 (Org-Level)} & \textbf{EIT-EMP-3 (Scope-Level)} \\
\hline
L0 & Individual & Systemic (langfristig) \\
L1 & Household & Strategic (jährlich) \\
L2 & Organization & Program (quartalsweise) \\
L3 & National & Operational (wöchentlich) \\
L4 & Global & Instant (< 24h) \\
\end{tabular}
\end{center}

\textbf{K* variiert nach Organizational Level (EIT-EMP-2):}
\begin{itemize}[nosep]
\item \textbf{L0 (Individual):} $K^* \approx 2-4$ (fokussiertes persönliches Portfolio)
\item \textbf{L1 (Household):} $K^* \approx 3-5$ (Koordination zwischen Partnern)
\item \textbf{L2 (Organization):} $K^* \approx 5-8$ (hierarchische Abstimmung)
\item \textbf{L3 (National):} $K^* \approx 4-6$ (Policy-Level, siehe Ch22)
\item \textbf{L4 (Global):} $K^* \approx 3-5$ (internationale Koordination, wenige aber tiefe Initiativen)
\end{itemize}

\textbf{Cross-Level Constraint:} $\gamma_{\text{L}_i, \text{L}_j} \geq 0$ --- Interventionen auf verschiedenen Org-Levels dürfen sich nicht gegenseitig unterminieren (vgl. Ch20 Theorem~\ref{thm:20-cross-level}).

\vspace{0.2cm}
\textit{Siehe: Chapter 11 §11.20 (EMERGE), Chapter 20 (Scope Levels L0-L4), Appendix UNMAPPED_IE}
\end{tcolorbox}

% =============================================================================
% INTUITION BOX
% =============================================================================

\begin{tcolorbox}[colback=yellow!5!white, colframe=yellow!75!black, title=Intuition: Anna Leads Organizational Change Across Levels]

\textbf{Anna} is Regional Health Director managing a major health initiative. She realizes quickly that success requires coordination \emph{across five different levels simultaneously}:

\begin{itemize}[nosep]
\item \textbf{Level 0 (Individual):} Individual doctors and nurses need to change how they work (adoption barrier: habit, skepticism)
\item \textbf{Level 1 (Household):} Patients come from families; if family doesn't support healthy behavior, patient effort fails (adoption barrier: family pressure)
\item \textbf{Level 2 (Org):} Hospital as organization needs to redesign workflows, budgets, training (adoption barrier: coordination costs)
\item \textbf{Level 3 (National):} Regional government sets funding levels and regulations (adoption barrier: political priorities shift)
\item \textbf{Level 4 (Global):} WHO guidelines and pharmaceutical companies' global pricing affect what's available locally (adoption barrier: external constraints)
\end{itemize}

Anna's mistake: She designed a brilliant Level 2 (organizational) portfolio and launched it. But ignoring Levels 1 (household), 3 (regional funding cuts), and 4 (medication shortages due to global supply) meant the initiative collapsed within 6 months.

\textbf{What Anna needed:} An integrated multi-level framework showing that 9C + Portfolio logic applies at \emph{every} level, and that spillovers between levels can sink even well-designed interventions.

This chapter shows how Anna (and you) can design and implement successfully across all five levels.

\end{tcolorbox}

% =============================================================================
% CENTRAL QUESTION BOX
% =============================================================================

\begin{tcolorbox}[colback=red!5!white, colframe=red!75!black, title=Central Question of This Chapter]

\textbf{How does the 9C + Portfolio framework scale from Individual to Global levels?}

\vspace{0.3cm}

More specifically:
\begin{itemize}
\item Does the 9C structure stay the same at each level? (YES --- framework is scale-invariant)
\item Do the parameters (L, d, γ, Ψ, Θ, A, W, φ) change at each level? (YES --- systematically)
\item What new challenges emerge at higher levels? (Aggregation, coordination, heterogeneity)
\item How do spillovers flow between levels? (Top-down policy → Bottom-up resistance)
\item When does bottom-up work better than top-down? (Depends on level heterogeneity)
\item How do we optimize implementation across \emph{all} levels simultaneously?
\end{itemize}

This chapter answers all six questions with theory + worked examples at each level.

\end{tcolorbox}

---

% =============================================================================
% SECTION 1: FRAMEWORK SCALE-INVARIANCE
% =============================================================================

\section{The Multi-Level Principle: 9C Works at Every Scale}
\label{sec:scale-invariance}

\textbf{Core insight:} The 9C diagnostic and portfolio optimization logic does \emph{not} depend on the level of analysis. Whether you're designing interventions for:

\begin{itemize}
\item A single person (Level 0: Individual)
\item A couple (Level 1: Household)
\item A department (Level 2: Organization)
\item A nation (Level 3: National)
\item Or the world (Level 4: Global)
\end{itemize}

the same structure applies:

\begin{equation}
\text{Diagnose } S_0 \to \text{ Design } T_i \to \text{ Modulate } (\alpha, \sigma) \to \text{ Optimize } P \to \text{ Execute } \& \text{ Learn}
\end{equation}

What changes:

\begin{enumerate}
\item \textbf{Unit of analysis:} Person → Household → Organization → Nation → Global
\item \textbf{Parameters:} How do L, d, γ, Ψ, A, W, φ \emph{quantitatively} differ?
\item \textbf{Constraints:} What coordination/aggregation problems emerge?
\item \textbf{Spillovers:} How do levels affect each other?
\end{enumerate}

The structure stays invariant; the specifics vary.

% =============================================================================
% SECTION 2: LEVEL 0 - INDIVIDUAL (FOUNDATION)
% =============================================================================

\section{Level 0: Individual --- The Foundation}
\label{sec:level0-individual}

\textbf{Chapters 16-21 cover Level 0 comprehensively.} Quick recap:

\begin{center}
\begin{tabular}{@{}lcc@{}}
\toprule
\textbf{9C Dimension} & \textbf{Unit} & \textbf{Range} \\
\midrule
L (Who) & Single person & Individual heterogeneity \\
d (What) & Utility components (FEPSDE) & 6 dimensions \\
γ (How) & 2-way interactions & Among 6 utilities \\
Ψ (Context) & 8 dimensions & Individual's environment \\
Θ (Parameters) & Sourced from literature & Intervention effectiveness \\
A (Awareness) & Subjective probability & 0\% to 100\% \\
W (Readiness) & Action threshold & 0 to 1 \\
φ (Phase) & BCJ stage (0-1) & PreAware → Maintenance \\
Scope & Individual level & Within-person decisions \\
\midrule
\end{tabular}
\end{center}

Example from Ch20-21: Individual with Type-2 diabetes deciding on behavior change. 9C diagnostic reveals: PreAware phase, Present-Biased segment, limited awareness of risk (A = 18\%).

Portfolio P₀ (Information + Feedback) increases A to 60\% in 3 months.

**Key Insight:** At Level 0, all variation is \emph{between} individuals (heterogeneous preferences, segments, contexts).

% =============================================================================
% SECTION 3: LEVEL 1 - HOUSEHOLD
% =============================================================================

\section{Level 1: Household --- Coupled Decisions}
\label{sec:level1-household}

\subsection{The Household 9C Framework}

When moving from individual to household, the 9C structure stays the same, but new dimensions emerge:

\begin{center}
\begin{tabular}{@{}lll@{}}
\toprule
\textbf{9C Dimension} & \textbf{Individual (L0)} & \textbf{Household (L1)} \\
\midrule
L (Who) & Single person & Multiple household members (partner, children) \\
d (What) & Individual utility & Household joint utility (U\_joint) \\
γ (How) & Individual's preferences interact & Preferences between household members \\
Ψ (Context) & Individual's environment & Shared household context \\
Θ (Parameters) & Individual effectiveness & Household coordination costs \\
A (Awareness) & Individual knowledge & \textbf{Asymmetric:} Different spouses know different things \\
W (Readiness) & Individual willingness & \textbf{Conflicting:} One spouse ready, other not \\
φ (Phase) & Individual's BCJ stage & \textbf{Mismatched:} Spouse A in Triggered, Spouse B in PreAware \\
Scope & Within-person & \textbf{Joint-level:} Decisions affect both \\
\midrule
\end{tabular}
\end{center}

\subsection{Household Complementarity (γ) Structure}

New type of γ: \textbf{Interpersonal complementarity} between spouses' decisions.

\textbf{Example:} Retirement savings decision

Partner A (female, 55, risk-averse): ``We should save aggressively now for retirement.'' (wants high $I_{\text{HOW}}$ commitment)

Partner B (male, 57, present-biased): ``We've waited this long, let's enjoy life now.'' (resists commitment)

\textbf{Complementarity analysis:}

\begin{quote}
If Partner A launches commitment strategy alone (high $I_{\text{HOW}}$), Partner B perceives this as ``controlling my money.'' Backlash: γ(A's commitment, B's willingness) = -0.4 (strong negative).

\textbf{Solution:} Joint commitment session where \emph{both} design the savings plan together. Reframes $I_{\text{HOW}}$ as ``our shared future.'' γ improves to +0.2 (weak positive).
\end{quote}

\subsection{Phase Mismatch Problem}

Real household challenge: Partners in \emph{different phases} of behavior change.

\textbf{Example:} Climate behavior

Partner A (φ = 0.7, Action phase): ``We need to transition to electric car NOW. I've been researching 6 months.''

Partner B (φ = 0.2, PreAware phase): ``But the charging stations don't exist here. I don't understand why we need to change.''

\textbf{Portfolio Design Problem:} What interventions work when partners are mismatched?

\textbf{Strategy:}
\begin{enumerate}
\item Bring Partner B to Triggered phase first: AWARE emphasis (information) + READY emphasis (test drive experience)
\item \emph{Then} joint commitment (HOW emphasis) when both are ready
\item Sequencing matters: Partner B needs AWARE$\rightarrow$READY before HOW
\end{enumerate}

Result: Phase mismatch resolved within 2-3 months; both partners agree on car purchase.

\subsection{Household Status Quo Diagnosis Template}

Using 9C framework for household:

\begin{center}
\small
\begin{tabular}{@{}llll@{}}
\toprule
\textbf{Dimension} & \textbf{Partner A} & \textbf{Partner B} & \textbf{Joint Issue} \\
\midrule
Phase & PreAware (0.3) & Triggered (0.6) & MISMATCH: Sequencing needed \\
Awareness & 25\% understand climate risk & 60\% aware & Asymmetric information \\
Readiness & Low (0.2) & Medium (0.5) & CONFLICT: Different willingness \\
Shared Context & Living situation, finances, family goals & Joint constraints on decisions \\
Interdependence & Both work; shared transport needs & γ between decisions \\
\midrule
\end{tabular}
\end{center}

% =============================================================================
% SECTION 4: LEVEL 2 - ORGANIZATION
% =============================================================================

\section{Level 2: Organization --- Hierarchical Coordination}
\label{sec:level2-organization}

\subsection{Organizational 9C Framework}

At organization level (department, company), 9C structure again invariant:

\begin{center}
\begin{tabular}{@{}lll@{}}
\toprule
\textbf{Dimension} & \textbf{Household (L1)} & \textbf{Organization (L2)} \\
\midrule
L (Who) & 2-3 household members & 50-5000+ employees \\
d (What) & Household utility & Organizational objectives + individual utilities \\
γ (How) & Spouse-spouse interactions & Manager-employee + peer interactions \\
Ψ (Context) & Household environment & Organizational culture + market conditions \\
Θ (Parameters) & Household coordination costs & Implementation/coordination costs \\
A (Awareness) & Spouse awareness mismatch & Department heads aware, frontline staff not \\
W (Readiness) & Partner disagreement & Management wants change, unions resist \\
φ (Phase) & Partner phase mismatch & Executives in Action, operations in PreAware \\
Scope & 2-person & Org-wide; cascading decisions \\
\midrule
\end{tabular}
\end{center}

\subsection{Organization as Nested Portfolio Problem}

Key insight: Organization = \emph{portfolio of portfolios}.

CEO designs org-level portfolio P(Org):
- Strategic initiatives ($I_{\text{WHO}}$, $I_{\text{WHEN}}$, $I_{\text{HOW}}$)
- Culture change programs ($I_{\text{WHO},o}$)

Department heads design dept-level portfolios P(Dept):
- Implementation of org initiatives + dept-specific tweaks
- Staffing, training ($I_{\text{AWARE}}$, $I_{\text{AWARE},\kappa}$)

Team leads design team-level portfolios P(Team):
- Daily execution
- Peer support, feedback ($I_{\text{AWARE},\kappa}$, $I_{\text{WHO},o}$)

\textbf{Coherence Problem:} What if dept portfolio conflicts with org portfolio?

\textbf{Example:} Manufacturing company restructuring

\textbf{Org-level portfolio (CEO):}
\begin{itemize}[nosep]
\item $I_{\text{WHEN}}$ (Choice architecture): ``New workflows; employees choose their own approach''
\item $I_{\text{WHO}}$ (Identity): ``You're now trusted professionals, not rule-followers''
\item $I_{\text{HOW}}$ (Commitment): ``Sign new 3-year employment contracts with profit-sharing''
\end{itemize}

\textbf{Problem:} Plant Manager overseeing restructuring designs different portfolio:
\begin{itemize}[nosep]
\item $I_{\text{WHEN},t}$ (Urgency): ``Minimize headcount by Q4. Execute fast, ask later.''
\item $I_{\text{WHAT},F}$ (Financial): ``Performance bonuses tied to cost-cutting.''
\item No choice architecture; no trust-building
\end{itemize}

\textbf{Conflict:} CEO says ``We're empowering you''; Plant Manager says ``Cut costs or lose your job.'' γ = -0.6 (strong negative).

\textbf{Result:} Unions file grievances; implementation collapses; turnover spikes.

\textbf{Solution:} Multi-level coherence review (Section 7 below).

\subsection{Manager as Intermediary}

Key role: \textbf{Manager translates org-level portfolio to team-level execution}.

Manager must:
1. Understand org-level portfolio (what CEO intended)
2. Translate to team context (what makes sense here?)
3. Maintain coherence (not contradict org-level message)
4. Add local context (this plant is different because...)

\textbf{Failure mode:} Manager doesn't understand org-level logic, waters down message, or reverses it.

% =============================================================================
% SECTION 5: LEVEL 3 - NATIONAL
% =============================================================================

\section{Level 3: National --- Regional Policy Implementation}
\label{sec:level3-national}

Chapters 22 + LLL + MMM cover national level in detail. Here we formalize it as Level 3:

\textbf{Level 3 = Regional/National Government implementing multi-domain portfolios}

Key differences from organizational:
\begin{itemize}
\item \textbf{Scope:} Population in millions (not hundreds)
\item \textbf{Heterogeneity:} Massive geographic/demographic variation in Ψ (urban vs. rural, wealthy vs. poor, educated vs. not)
\item \textbf{Multiple domains:} Health, pension, environment, labor all simultaneously
\item \textbf{Cross-domain spillovers:} Must manage γ between domains (see Ch22)
\item \textbf{Decentralization:} Cities/regions implement differently
\item \textbf{Politics:} Opposition parties, interest groups, media can undermine
\end{itemize}

\subsection{Regional Heterogeneity Problem}

Major challenge at national level: \textbf{Same portfolio, different Ψ, different effectiveness}.

\textbf{Example:} Energy efficiency program (from Appendix UNMAPPED_LLL)

\textbf{Urban Region (high-income, educated):}
- High Ψ\_K (information quality)
- High Ψ\_F (factor flexibility; can afford upgrades)
- Awareness already 40\%
- Energy efficiency program effectiveness: 14\%

\textbf{Rural Region (low-income, low education):}
- Low Ψ\_K (information quality; unreliable internet)
- Low Ψ\_F (can't afford home upgrades)
- Awareness only 8\%
- Same program effectiveness: 4\% (massive drop)

\textbf{National Government Problem:} Design portfolio for \emph{national} level that works in \emph{both} regions.

\textbf{Solution:} Segment-and-target approach

- Urban region: Keep program as-is
- Rural region: Pre-invest in information infrastructure ($I_{\text{AWARE}}$ expansion) + government co-funding ($I_{\text{WHAT},F}$) before efficiency program

Result: Both regions reach ~12\% effectiveness (higher for rural than baseline 4%, lower for urban than 14\%, but acceptable compromise).

% =============================================================================
% SECTION 6: LEVEL 4 - GLOBAL
% =============================================================================

\section{Level 4: Global --- International Coordination}
\label{sec:level4-global}

\subsection{The Global Coordination Challenge}

Global level = \textbf{Multi-nation coordination without global government}

\textbf{Example:} Climate Agreement

Imagine negotiating a climate policy portfolio that requires behavior change in 195 countries with:
- Radically different Ψ contexts (carbon-intensive economies vs. renewables)
- Conflicting interests (developed vs. developing nations)
- No enforcement mechanism (unlike national government)
- Long time horizons (20-50 years)
- Massive distributional conflicts (who pays for transition?)

\subsection{Global 9C Framework}

\begin{center}
\begin{tabular}{@{}lll@{}}
\toprule
\textbf{Dimension} & \textbf{National (L3)} & \textbf{Global (L4)} \\
\midrule
L (Who) & Population of nation & All humans (or subset: all nations) \\
d (What) & National objectives & Global objectives (climate stability) + national interests \\
γ (How) & Cross-domain spillovers & Countries' interests interact \\
Ψ (Context) & Regional variation & Massive variation: Arctic vs. Sahara vs. Island \\
Θ (Parameters) & Implementation in regions & Implementation in countries; sovereignty \\
A (Awareness) & National public opinion & Global public opinion (highly varied) \\
W (Readiness) & National political will & National commitments to treaty \\
φ (Phase) & Policy phase nationally & Country heterogeneity in phase \\
Scope & National & International; no enforcement \\
\midrule
\end{tabular}
\end{center}

\subsection{Global Spillover Example: Carbon Pricing}

Suppose we design global portfolio: ``Each nation implements carbon tax starting 2025.''

\textbf{Spillover Problem 1: Trade Effects}

Country A (strong climate commitment, high carbon tax): Manufacturing becomes expensive; jobs move to Country B

Country B (weaker commitment, low carbon tax): Gets manufacturing jobs BUT increases total global emissions (net negative)

Complementarity: γ(A's carbon tax, B's carbon tax) = NEGATIVE if not coordinated

Solution: Border adjustment taxes; carbon tariffs; negotiated phase-in

\textbf{Spillover Problem 2: Political Coalitions}

Within Country A: Progressive voters support carbon tax; Manufacturing workers oppose

Within Country B: Low-income voters oppose (energy bills rise); Industrialists support (tariff protection)

Complementarity: What helps Country A politically hurts Country A environmentally if wrong coalition wins power

\subsection{Bottom-Up vs Top-Down at Global Level}

\textbf{Top-down (Treaty-based):}
- Negotiated agreement signed by governments
- Binding commitments (theoretically)
- Problem: Enforcement weak; countries defect if conditions change

\textbf{Bottom-up (Market-based):}
- Individual companies commit to carbon neutrality
- Supply chain pressure drives behavior
- Cities/regions implement unilaterally
- Problem: Fragmented; may not reach scale needed

\textbf{Hybrid (increasingly common):}
- Treaty sets goals; implementation varies by context
- Global framework + local flexibility
- Example: Paris Climate Agreement (goals global, implementation national)

% =============================================================================
% SECTION 7: CROSS-LEVEL SPILLOVERS
% =============================================================================

\section{Cross-Level Spillovers and Constraints}
\label{sec:cross-level-spillovers}

\subsection{Top-Down Causation: How National Policy Affects Household}

\textbf{Example:} Pension reform (Level 3) affects household savings decisions (Level 1)

\textbf{National policy:} ``Retirement age increases from 65 to 67; benefits reduced 20\%'' (Level 3 portfolio: $I_{\text{WHEN},t}$ temporal urgency + $I_{\text{WHEN}}$ choice architecture)

\textbf{Household response:} Different households respond differently depending on context

\begin{center}
\begin{tabular}{@{}lll@{}}
\toprule
\textbf{Household Type} & \textbf{Response} & \textbf{Effectiveness} \\
\midrule
High-income couple (both employed) & Increase retirement savings ($I_{\text{HOW}}$ commitment works) & +25\% savings rate \\
Middle-income couple (one job) & Stress about not having enough; save anxiously & +8\% savings (insufficient) \\
Low-income couple (both precarious work) & Panic; little ability to save more & +0\% (can't afford to change) \\
\midrule
\end{tabular}
\end{center}

\textbf{National Problem:} Policy assumed middle-income response (+8\%) but needed +15\% nationally to make reform work. Low-income households couldn't comply.

\textbf{Solution:} Subsidized retirement account for low-income ($I_{\text{WHAT},F}$ + $I_{\text{WHEN}}$); increased benefits for those who can't work to 67 ($I_{\text{WHEN}}$ flexibility).

Result: Aggregate response: +12\% (closer to target).

\subsection{Bottom-Up Pressure: How Individual/Household Behavior Constrains National Policy}

\textbf{Reverse direction:} Individual choices aggregate up to make national policy politically infeasible.

\textbf{Example:} Energy efficiency policy (Level 3) requires household behavior change (Level 1)

National government announces: ``Energy tax +20\% to fund renewable transition. Efficiency programs offset costs.''

\textbf{Expected:} Households reduce consumption via efficiency programs (portfolio: $I_{\text{AWARE}}$ information + $I_{\text{AWARE},\kappa}$ feedback + $I_{\text{WHAT},F}$ incentive)

\textbf{Reality:} 40\% of households reduce consumption (early adopters); 60\% don't (resistant or unable)

Political pressure:
- Opposition parties mobilize the 60\%: ``Government is hurting the poor!''
- Protests grow
- Government backtracks before policy achieves goals

\textbf{Cross-level spillover:} Individual non-compliance → Political failure → National policy collapses

\textbf{Solution:} Better sequencing ($I_{\text{AWARE}}$ information → $I_{\text{AWARE},\kappa}$ feedback → THEN $I_{\text{WHAT},F}$ incentive, not all at once); targeted support for vulnerable groups; manage political narrative in parallel.

\subsection{Organizational Cascade: How National Policy Affects Organization Affects Team}

Three-level cascade:

\begin{center}
\begin{tabular}{@{}llll@{}}
\toprule
\textbf{Level} & \textbf{Actor} & \textbf{Message} & \textbf{Actual Effect} \\
\midrule
Level 3 & National Govt & ``Reduce energy 20\%'' & 20\% target \\
↓ & & & \\
Level 2 & Company CEO & ``Our company will lead; 30\% reduction'' & 30\% target (!)  \\
& & (Misinterprets; overcommits) & \\
↓ & & & \\
Level 1 & Plant Manager & ``Shutdown energy-intensive line; fire 200 people'' & Layoffs, union strike \\
& & (Sees only cost-cutting) & 0\% reduction \\
\midrule
\end{tabular}
\end{center}

Cascade failure: Message transformed at each level; final implementation is opposite of intent.

% =============================================================================
% SECTION 8: BOTTOM-UP VS TOP-DOWN IMPLEMENTATION
% =============================================================================

\section{When Does Bottom-Up Work Better Than Top-Down?}
\label{sec:bottom-up-top-down}

\subsection{Bottom-Up Advantages}

\begin{itemize}
\item \textbf{Local adaptation:} Teams understand their context; can customize portfolio to local Ψ
\item \textbf{Buy-in:} People more willing to comply when they co-designed the solution
\item \textbf{Information:} Local teams have tacit knowledge that central planners miss
\item \textbf{Speed:} Don't wait for central approval; pilot locally, scale what works
\end{itemize}

\textbf{Best for:} Implementation of well-understood goals (everyone agrees on what, not how)

\textbf{Example:} ``Reduce diabetes prevalence 20\%'' — local health departments know their populations; let them design portfolios.

\subsection{Top-Down Advantages}

\begin{itemize}
\item \textbf{Coordination:} When activities must be synchronized (all cutover on same date for system migration)
\item \textbf{Equity:} Ensures minimum standards even in disadvantaged areas (wouldn't happen naturally bottom-up)
\item \textbf{Network effects:} Some benefits only appear when everyone does it (financial system stability)
\item \textbf{Speed for complex coordination:} Faster than negotiating bottom-up consensus
\end{itemize}

\textbf{Best for:} Big system changes, equity requirements, network effects

\textbf{Example:} ``All hospitals switch to new electronic health records by June 1'' — needs top-down mandate; bottom-up would create incompatible systems.

\subsection{Hybrid Model (Most Effective)}

\textbf{Top-down:} Set goals, budgets, minimum standards

\textbf{Bottom-up:} Design and execute implementation

\textbf{Example:} GDPR (European data protection regulation)

- Top-down: ``All companies must protect personal data; right to be forgotten; etc.'' (EU mandate, non-negotiable)
- Bottom-up: ``Each company decides how to implement; hire DPO (data protection officer); audit your processes'' (company-level design)

Result: Universal compliance + company-specific solutions = better than pure top-down OR pure bottom-up

% =============================================================================
% SECTION 9: INTEGRATED MULTI-LEVEL CASE STUDY
% =============================================================================

\section{Integrated Case Study: Health System Reform Across All Levels}
\label{sec:case-study-health}

\subsection{The Challenge}

Country: Medium-income nation with fragmented health system. Goals:
- Increase preventive care (reduce expensive acute care)
- Improve primary care quality
- Ensure equity (rural areas underserved)
- Fiscal sustainability (10\% budget reduction needed)

Portfolio must work across:

**Level 4 (Global):** WHO guidance, pharmaceutical pricing, medical device availability

**Level 3 (National):** Health ministry must set policy and budget

**Level 2 (Organization):** Public health system + private providers + NGOs

**Level 1 (Household):** Patient families must adopt preventive behaviors

**Level 0 (Individual):** Doctors, nurses, patients must change daily work

\subsection{Level 0-3 Portfolio Design}

\textbf{Level 0 (Individual Doctor/Patient):}
- Doctor: $I_{\text{AWARE}}$ (training on prevention protocols) + $I_{\text{AWARE},\kappa}$ (feedback on patient outcomes)
- Patient: $I_{\text{AWARE}}$ (health education) + $I_{\text{AWARE},\kappa}$ (feedback from wearables) + $I_{\text{WHO},o}$ (peer support for behavior change)

\textbf{Level 1 (Household):}
- Family health counseling ($I_{\text{WHO},o}$: peer support with other families)
- Shared health savings account ($I_{\text{HOW}}$: commitment to joint budget for prevention)
- Partner accountability (if one spouse skips appointments, other reminds)

\textbf{Level 2 (Organization):}
- Health facility redesign ($I_{\text{WHEN}}$: walk-ins for preventive care; no appointment needed)
- Staff empowerment ($I_{\text{WHO}}$: nurses now do initial preventive care, not just admin)
- Quality reporting ($I_{\text{AWARE},\kappa}$: transparency about health outcomes by facility)

\textbf{Level 3 (National):}
- Payment reform ($I_{\text{WHAT},F}$: clinics paid for prevention, not just treatment)
- Rural incentive ($I_{\text{WHAT},F}$: bonuses for clinics in underserved areas who hit targets)
- Regulation ($I_{\text{WHEN},t}$: mandatory preventive screenings covered by insurance)

\subsection{Cross-Level Spillovers Managed}

\textbf{Positive spillovers deliberately created:}

Level 3 → Level 2: Payment reform ($I_{\text{WHAT},F}$) makes investment in facilities ($I_{\text{WHEN}}$) attractive
Level 2 → Level 1: Quality transparency ($I_{\text{AWARE},\kappa}$) motivates household engagement ($I_{\text{WHO},o}$)
Level 1 → Level 0: Family support ($I_{\text{WHO},o}$) enables doctor compliance ($I_{\text{AWARE}}$ training more effective)

\textbf{Negative spillovers actively prevented:}

Risk: Level 3 budget cuts (fiscal austerity) conflict with Level 2 investment needs
Solution: Payment reform stays protected even during budget cuts

Risk: Level 1 Household conflict: ``Why should we spend on prevention when we can barely afford food?''
Solution: Subsidized prevention; target prevention first to high-risk families

\subsection{Implementation Timeline}

\textbf{Year 1: Level 0 + 2 focus}
- Doctor training ($I_{\text{AWARE}}$)
- Facility redesign ($I_{\text{WHEN}}$)
- Expected effect: +8\% preventive visits

\textbf{Year 2: Level 1 + 3 addition}
- Household engagement program ($I_{\text{WHO},o}$)
- Payment reform ($I_{\text{WHAT},F}$)
- Rural incentives ($I_{\text{WHAT},F}$)
- Expected effect: +15\% preventive visits (cumulative)

\textbf{Year 3: Full integration + Level 4 alignment}
- All levels coordinated
- Global guidelines fully integrated
- Expected effect: +25\% preventive visits
- Budget savings: 10\% achieved through prevention

\subsection{Success Indicators by Level}

\begin{center}
\small
\begin{tabular}{@{}lccc@{}}
\toprule
\textbf{Level} & \textbf{Metric} & \textbf{Y1 Target} & \textbf{Y3 Target} \\
\midrule
L0 (Individual) & Doctors completing training & 80\% & 95\% \\
L1 (Household) & Families in prevention program & 30\% & 70\% \\
L2 (Organization) & Health facilities with new protocols & 50\% & 100\% \\
L3 (National) & Population receiving preventive care & 15\% & 45\% \\
L4 (Global) & WHO guidelines implemented & 60\% & 95\% \\
\midrule
\end{tabular}
\end{center}

Key finding: Y3 success required progress at \emph{all five levels} simultaneously. If any level got stuck, whole system slowed.

---

% =============================================================================
% READING PATHS
% =============================================================================

\section*{Reading Paths and Integration}

\begin{tcolorbox}[colback=yellow!5!white, colframe=yellow!75!black, title=Reading Path: What Comes Next?]

\textbf{You have now completed:}
\begin{itemize}[nosep]
\item Chapters 16-23: Individual diagnosis → multi-level scaling
\item EMERGE Framework (Ch11): $K^*$ variiert nach Org-Level (L0: 2-4, L2: 5-8, L4: 3-5)
\end{itemize}

\textbf{Next step (Chapter 24 - Emergent Life Journeys):}

Go to Chapter 24 to understand how the 7 life journeys \textbf{EMERGE} from domain-specific parameters over 30-50 year lifespans. This chapter formalizes why intervention timing differs so dramatically per domain:
\begin{itemize}[nosep]
\item Health converges in 5-10 months → maintenance after year 1
\item Money converges SLOWLY (30+ years) → continuous intervention required
\item Career is episodic (per-role learning) → intensive support at transitions
\end{itemize}

\textbf{Chapter 24 Resources:}
\begin{itemize}[nosep]
\item Appendix UNMAPPED_RRR (METHOD-LIFETIME) - 30-year templates per domain
\item Appendix UNMAPPED_OOO (FORMAL-SYSTEM-DYNAMICS) - Complete ODE analysis
\item Appendix UNMAPPED_PPP (METHOD-AGENT-BASED) - Agent-based simulation code
\end{itemize}

\end{tcolorbox}

\subsection{Formal Details}
\label{sec:ch23-formal}

The complete formal treatment appears in Appendix UNMAPPED_NNN (METHOD-MULTI-LEVEL). Key results:

\begin{itemize}[nosep]
    \item \textbf{Level Aggregation:} Micro-to-macro aggregation functions for behavioral outcomes
    \item \textbf{Cross-Level γ:} Complementarity matrices that capture spillovers between levels
    \item \textbf{Cascade Dynamics:} How policy signals transform as they pass through organizational hierarchy
    \item \textbf{Coherence Conditions:} Mathematical criteria for portfolio alignment across levels
    \item \textbf{Optimal Hybrid:} When to use top-down vs bottom-up implementation strategies
\end{itemize}

\subsection{What Comes Next}
\label{sec:ch23-what-comes-next}

\begin{enumerate}
    \item \textbf{Chapter 11 §11.20 (EMERGE):} Algorithmus für $K^*$ auf jedem Organizational Level
    \item \textbf{Chapter 24 (Emergent Life Journeys):} Domain-spezifische 30-50 Jahre Trajektorien
    \item \textbf{Appendix UNMAPPED_NNN:} Full mathematical treatment of multi-level coordination
    \item \textbf{Appendix UNMAPPED_IE (EMERGE-SPEC):} Scope-Level vs. Org-Level Unterscheidung
\end{enumerate}

%======================================================================
% SUMMARY SECTION
%======================================================================

\section{Summary}
\label{sec:ch23-summary}

\begin{tcolorbox}[colback=blue!5!white, colframe=blue!75!black,
    title={\textbf{Chapter 23 Summary: Multi-Level Implementation}}]

\textbf{Core Insight:} The 9C + Portfolio framework is scale-invariant---the same diagnostic and optimization logic applies at Individual (L0), Household (L1), Organization (L2), National (L3), and Global (L4) levels.

\textbf{Key Results:}
\begin{enumerate}[nosep]
    \item \textbf{Structure Invariance:} 9C dimensions exist at every organizational level (L0-L4)
    \item \textbf{K* by Level (EIT-EMP-2):} Optimale Portfolio-Größe variiert: L0 (2-4) → L2 (5-8) → L4 (3-5)
    \item \textbf{Cross-Level Constraint:} $\gamma_{\text{L}_i, \text{L}_j} \geq 0$ (keine gegenseitige Unterminierung)
    \item \textbf{Two Level-Concepts:} Org-Levels (Ch23) $\neq$ Scope-Levels (EIT-EMP-3)
    \item \textbf{Hybrid Optimality:} Top-down goals + bottom-up execution outperforms pure approaches
\end{enumerate}

\textbf{Transition:} Chapter 24 examines how domain-specific parameters create emergent life journeys over 30-50 year timescales.

\end{tcolorbox}

\begin{center}
\textbf{One Framework, Five Levels:} \\
\vspace{0.3cm}

\textbf{Level 0: Individual} (Foundation) \\
Chapters 16-21 \\
↓ \\
\textbf{Level 1: Household} (Coupled decisions, phase mismatch) \\
New: Ch23 Section 3 \\
↓ \\
\textbf{Level 2: Organization} (Hierarchical coordination) \\
New: Ch23 Section 4 \\
↓ \\
\textbf{Level 3: National} (Chapters 22, LLL, MMM) \\
Formalized: Ch23 Section 5 \\
↓ \\
\textbf{Level 4: Global} (International coordination) \\
New: Ch23 Section 6 \\
\vspace{0.5cm}

\textbf{Integrated Implementation:} \\
Cross-Level Spillovers + Bottom-Up/Top-Down Hybrid \\
(Ch23 Sections 7-9)
\end{center}

%======================================================================
% END OF CHAPTER 23
%======================================================================

\newpage
